% page 207 of the facsimile (image p-206 of the scan)
% block 170.2 x 242.0 mm ; left margin 31.8 mm
\pgc{10.160mm}{0.000mm}{
\l{}
\l{}
\l{}
\l{\cel{58}199}
\l{}
\l{\sou{graco}.\cel{2}(\sou{teol.})\cel{2}Sokurso supernatura quan Deo donas a la homo por}
\l{\cel{3}helpar lu salvar su. - EFIRS.}
\l{}
\l{\sou{gracio}. Agreablajo, plezivajo, en persono, en la kozi. - DEFIRS.}
\l{}
\l{\sou{gracila}.\cel{2}Gracioze tenua ed alta-statura. - EFIS.}
\l{}
\l{\sou{grado}.- I. La parto di eskalero, di skalo, sur qua onu pozas la}
\l{\cel{3}pedo por acensar o decensar. - II. Singla de la mediataji qui}
\l{\cel{3}duktas sucede de ca a ta estado. - III. En serio de termini di}
\l{\cel{3}qui la importo, la valoro, kreskas o deskreskas, plaso di un}
\l{\cel{3}ek la termini relate la ceteri. - IV. Singla ek la dividuri}
\l{\cel{3}inter-egala di serio markoza. - DEFIRS.}
\l{}
\l{\sou{gradiento}. (\sou{meteorol.}) Difero di la preso atmosferala, evaluata per}
\l{\cel{3}milimetri e per grado geografiala, inter punto e la centro maxim}
\l{\cel{3}apuda di ciklono o di anticiklono.}
\l{}
\l{\sou{grafika}. I. Qua trasas per desegno. - II. Qualeso di la trasuro da}
\l{\cel{3}aparato enrejistrala. - III. Qua relatas la reprezento di la}
\l{\cel{3}son-uni vocala per la skribado. - DEFIRS.}
\l{}
\l{\sou{grafism}o.I. Maniero skribar la vorti di linguo. - II. Skribo-stilo}
\l{\cel{3}di individuo, e qua (segun-dice) informas pri lua karaktero.}
\l{}
\l{\sou{grafito}. (\sou{mineral.}) Substanco (\sou{ferkarburo}) de qua onu facas la kra-}
\l{\cel{3}yoni ordinara. - DEFIRS.}
\l{}
\l{\sou{grafologio}. Studiado di la karaktero di la homi segun lia skribo-}
\l{\cel{3}stilo, di lia grafismo. - DEFIRS.}
\l{}
\l{\sou{grafometro.}\cel{2}Mi-cirklo kun dioptro ed alidado, por mezurar la}
\l{\cel{3}anguli lor agromezuro, mapo-relevo.}
\l{}
\l{gramo.\cel{2}Ponder-unajo di la metro-sistemo, to quon pezas un centi-}
\l{\cel{3}metro kuba de aquo distilita, ye 4 gradi Celsius. - DEFIRS.}
\l{}
\l{gramatiko.\cel{2}Cienco pri la reguli di linguo. - DEFIRS.}
\l{}
\l{gramino.\cel{2}(\sou{bot.)}\cel{2}Familio de planti monokotiledona, qui havas}
\l{\cel{3}kom stipo palio kava, nodoza, kun flori dispozita spike.}
\l{}
\l{grano.\cel{2}I.Singla de la frukti kontenita en la spiko di la cerealo.}
\l{\cel{3}- II. Frukto o semino globatra di ula planti (vito, grozeliero).}
\l{\cel{3}- III. Asperajo granatra di surfaco (petro, stofo, e c.). -}
\l{\cel{3}IV. Karaktero granoza di ca surfaco. - FISDE.}
\l{}
\l{granario. La parto interna maxim supra di domo, uzata kom depozeyo}
\l{\cel{3}desembarasala. - FIS.}
\l{}
\l{granato. Lapido, vino-reda, qua povas striizar quarco. - DEFIRS.}
}
